\documentclass[main.tex]{subfiles}

\begin{document}

\section{INTRODUCTION}

The Motor Control Box is the advanced module of the PLC Lab Trainer System, building on the foundational concepts introduced by the Heater Control Box.

The purpose of this project is to develop a basic motor control board specifically designed to drive small, low-cost DC brushed motors using industrial signals. This innovative board is capable of accepting 0-10V inputs from a Programmable Logic Controller (PLC), allowing precise control over the commanded motor power. Additionally, it supports 24V Pulse Width Modulation (PWM) control signals directly for both forward and reverse motor operations, enhancing its versatility and applicability for various learning applications.

The motor control board is equipped with feedback capabilities, providing 0-10V outputs to the PLC for motor speed and position and a digital signal for motor direction. This allows a diverse range of process control learning opportunities. Furthermore, the board includes a position zeroing input from the PLC, enabling the control board to reset the motor position feedback to zero when desired to allow a homing routine to be used to establish absolute position control.

Overall, this project will provide a versatile motor controller that will facilitate a diverse range of learning opportunities.

\begin{figure}[H]
    \centering
    \includegraphics[width=\textwidth]{images/MotorBoxConcept.png}
    \caption{Motor Control Box system diagram. The PLC issues a 0--10\,V power command, 24\,V PWM forward/reverse signals, and a position-zero input. A Raspberry Pi Pico converts those industrial signals to logic-level commands for the H-Bridge driver, which powers the DC brushed motor. The motor's quadrature encoder is decoded by the Pico's PIO and returned to the PLC as 0--10\,V analog speed and position plus a digital direction signal.}
    \label{fig:motorbox-concept}
\end{figure}



\end{document}